\documentclass[11pt,a4paper]{article}

%=========================================================
% PACKAGES
%=========================================================
\usepackage[utf8]{inputenc}
\usepackage[T1]{fontenc}
\usepackage[french]{babel}
\usepackage[a4paper,margin=2.2cm]{geometry}
\usepackage{lmodern}
\usepackage{amsmath,amssymb,amsthm,mathtools}
\usepackage{fancyhdr}
\usepackage{lastpage}
\usepackage{enumitem}
\usepackage{array,tabularx}
\usepackage{tikz}
\usepackage[most]{tcolorbox}
\usepackage{hyperref}
\usepackage{xcolor}

%=========================================================
% MISE EN PAGE
%=========================================================
\setlength{\headheight}{14pt}
\pagestyle{fancy}
\fancyhf{}
\fancyhead[L]{Mathématiques CPGE}
\fancyhead[C]{Polynômes}
\fancyhead[R]{Page \thepage/\pageref{LastPage}}
\fancyfoot[C]{Cours complet -- Algèbre}

\hypersetup{
    colorlinks=true,
    linkcolor=blue!50!black,
    urlcolor=blue!50!black
}

\setlist[itemize]{label=\(\triangleright\)}
\setlist[enumerate]{label=\textbf{\arabic*.}}

%=========================================================
% COULEURS ET BOÎTES
%=========================================================
\definecolor{bleufonce}{RGB}{25,65,130}
\definecolor{bleuclair}{RGB}{235,243,255}
\definecolor{vertfonce}{RGB}{30,120,70}
\definecolor{vertclair}{RGB}{235,250,240}
\definecolor{rougefonce}{RGB}{170,40,40}
\definecolor{rougeclair}{RGB}{255,238,238}
\definecolor{orangefonce}{RGB}{190,110,20}
\definecolor{orangeclair}{RGB}{255,246,230}
\definecolor{violetfonce}{RGB}{95,55,145}
\definecolor{violetclair}{RGB}{246,241,255}

\tcbset{
    enhanced,
    sharp corners,
    boxrule=0.7pt,
    before skip=8pt,
    after skip=8pt
}

\newtcolorbox{definitionbox}[1][]{
    colback=bleuclair,
    colframe=bleufonce,
    title=\textbf{Définition #1}
}

\newtcolorbox{theorembox}[1][]{
    colback=vertclair,
    colframe=vertfonce,
    title=\textbf{Théorème / Propriété #1}
}

\newtcolorbox{methodbox}[1][]{
    colback=orangeclair,
    colframe=orangefonce,
    title=\textbf{Méthode #1}
}

\newtcolorbox{retenirbox}[1][]{
    colback=violetclair,
    colframe=violetfonce,
    title=\textbf{À retenir #1}
}

\newtcolorbox{erreurbox}[1][]{
    colback=rougeclair,
    colframe=rougefonce,
    title=\textbf{Erreur classique #1}
}

\newtcolorbox{exercicebox}[1][]{
    colback=white,
    colframe=black!60,
    title=\textbf{Exercice #1}
}

\newtcolorbox{correctionbox}[1][]{
    colback=gray!5,
    colframe=black!65,
    title=\textbf{Correction #1}
}

%=========================================================
% COMMANDES
%=========================================================
\newcommand{\K}{\mathbb K}
\newcommand{\R}{\mathbb R}
\newcommand{\C}{\mathbb C}
\newcommand{\N}{\mathbb N}
\newcommand{\Q}{\mathbb Q}
\newcommand{\Z}{\mathbb Z}
\newcommand{\KX}{\K[X]}
\newcommand{\RX}{\R[X]}
\newcommand{\CX}{\C[X]}
\newcommand{\KnX}{\K_n[X]}
\newcommand{\degp}{\operatorname{deg}}
\newcommand{\pgcd}{\operatorname{pgcd}}
\newcommand{\ppcm}{\operatorname{ppcm}}
\newcommand{\Vect}{\operatorname{Vect}}
\newcommand{\Ker}{\operatorname{Ker}}
\newcommand{\Id}{\operatorname{Id}}
\newcommand{\Mat}{\operatorname{Mat}}

\newtheorem{prop}{Proposition}[section]
\newtheorem{lemme}[prop]{Lemme}
\newtheorem{corollaire}[prop]{Corollaire}

%=========================================================
% DOCUMENT
%=========================================================
\begin{document}

\begin{center}
    {\Huge \textbf{Polynômes}}\\[2mm]
    {\Large Structure de \(\K[X]\), racines, factorisation, arithmétique et interpolation}\\[2mm]
    {\large Cours complet pour classes préparatoires scientifiques}\\[2mm]
    \rule{0.88\linewidth}{0.5pt}
\end{center}

\vspace{0.4cm}

\tableofcontents

\newpage

%=========================================================
\section*{Introduction générale}
\addcontentsline{toc}{section}{Introduction générale}

Les polynômes sont des objets déjà rencontrés au lycée sous la forme de fonctions :
\[
x\longmapsto ax^2+bx+c,
\qquad
x\longmapsto x^3-2x+1,
\]
mais en classe préparatoire, ils deviennent des objets algébriques à part entière.

Un polynôme n'est pas seulement une fonction.  
C'est d'abord une expression formelle :
\[
P(X)=a_0+a_1X+\cdots+a_nX^n,
\]
où \(X\) est une indéterminée et où les coefficients \(a_k\) appartiennent à un corps \(\K\).

\begin{retenirbox}
Un polynôme est un objet algébrique formel.  
Il peut ensuite être évalué en un scalaire pour donner une fonction polynomiale, mais il ne faut pas confondre le polynôme et la fonction qu'il définit.
\end{retenirbox}

Les polynômes interviennent partout en mathématiques :
\begin{itemize}
    \item en algèbre : divisibilité, factorisation, PGCD, Bézout ;
    \item en analyse : dérivation, Taylor, interpolation ;
    \item en algèbre linéaire : polynômes caractéristiques, polynômes annulateurs, diagonalisation ;
    \item en arithmétique : analogie forte entre \(\Z\) et \(\K[X]\).
\end{itemize}

Le but de ce chapitre est de construire progressivement la théorie de \(\K[X]\), en insistant sur les démonstrations et sur les méthodes pratiques.

%=========================================================
\section{Qu'est-ce qu'un polynôme ?}

\subsection{L'indéterminée \(X\)}

Avant de définir rigoureusement un polynôme, il faut comprendre le rôle de la lettre \(X\).

Dans l'expression :
\[
P(X)=3X^4-2X+1,
\]
le symbole \(X\) n'est pas un nombre.  
C'est une \textbf{indéterminée}.  
Cela signifie que l'on manipule formellement des puissances :
\[
1,\ X,\ X^2,\ X^3,\dots
\]
avec des coefficients dans un corps \(\K\).

\begin{erreurbox}
La lettre \(X\) dans \(\K[X]\) n'est pas une variable réelle au sens analytique.  
C'est une indéterminée formelle.  
Le polynôme existe avant toute évaluation en un nombre.
\end{erreurbox}

\subsection{Définition formelle}

\begin{definitionbox}[Polynôme à coefficients dans \(\K\)]
Soit \(\K\) un corps.  
Un polynôme à coefficients dans \(\K\) est une expression formelle :
\[
P(X)=a_0+a_1X+\cdots+a_nX^n,
\]
où :
\[
n\in\N,
\qquad
a_0,\dots,a_n\in\K.
\]
L'ensemble des polynômes à coefficients dans \(\K\) est noté :
\[
\K[X].
\]
\end{definitionbox}

On peut aussi écrire :
\[
P(X)=\sum_{k=0}^n a_kX^k.
\]

Les scalaires \(a_k\) sont appelés les \textbf{coefficients} du polynôme.

\begin{retenirbox}
Un polynôme est une somme finie de puissances de \(X\) pondérées par des coefficients.
\[
P(X)=\sum_{k=0}^n a_kX^k.
\]
\end{retenirbox}

\subsection{Égalité de deux polynômes}

\begin{definitionbox}[Égalité de deux polynômes]
Deux polynômes :
\[
P(X)=\sum_{k=0}^n a_kX^k,
\qquad
Q(X)=\sum_{k=0}^m b_kX^k
\]
sont égaux si et seulement si leurs coefficients correspondants sont égaux :
\[
\forall k\in\N,\qquad a_k=b_k.
\]
On convient que les coefficients absents valent \(0\).
\end{definitionbox}

\begin{exercicebox}[Compréhension]
Les deux écritures suivantes définissent-elles le même polynôme ?
\[
P(X)=X^2+1,
\qquad
Q(X)=0X^4+0X^3+X^2+0X+1.
\]
\end{exercicebox}

\begin{correctionbox}
Oui. Les coefficients de \(X^4\), \(X^3\) et \(X\) dans \(Q\) sont nuls.  
Donc les deux polynômes ont exactement les mêmes coefficients :
\[
a_2=1,\quad a_0=1,
\]
et tous les autres coefficients sont nuls.
\[
P=Q.
\]
\end{correctionbox}

\begin{erreurbox}
Il faut distinguer :
\begin{itemize}
    \item l'égalité formelle de polynômes ;
    \item l'égalité de fonctions polynomiales.
\end{itemize}

Sur un corps infini, comme \(\R\) ou \(\C\), deux polynômes qui définissent la même fonction sont égaux.  
Sur un corps fini, ce n'est pas toujours vrai.
\end{erreurbox}

%=========================================================
\section{Structure algébrique de \(\K[X]\)}

\subsection{Addition et multiplication}

Soient :
\[
P(X)=\sum_{k=0}^n a_kX^k,
\qquad
Q(X)=\sum_{k=0}^m b_kX^k.
\]

On définit :
\[
(P+Q)(X)=\sum_{k\geq0}(a_k+b_k)X^k.
\]

Le produit est défini par :
\[
(PQ)(X)=\sum_{k\geq0} c_kX^k,
\]
où :
\[
c_k=\sum_{i+j=k}a_ib_j.
\]

\begin{theorembox}[Structure d'anneau]
L'ensemble \(\K[X]\), muni de l'addition et de la multiplication usuelles des polynômes, est un anneau commutatif unitaire.
\end{theorembox}

\begin{proof}
On vérifie les axiomes directement à partir des coefficients.

\begin{itemize}
    \item L'addition est associative et commutative car l'addition dans \(\K\) l'est.
    \item Le polynôme nul est l'élément neutre pour l'addition :
    \[
    0_{\K[X]}=0.
    \]
    \item Tout polynôme \(P\) admet un opposé \(-P\), obtenu en changeant le signe de tous ses coefficients.
    \item La multiplication est associative, commutative et distributive par rapport à l'addition, car ces propriétés sont vraies dans \(\K\).
    \item Le polynôme constant \(1\) est l'élément neutre multiplicatif.
\end{itemize}

Donc \(\K[X]\) est un anneau commutatif unitaire.
\end{proof}

\begin{theorembox}[Structure d'espace vectoriel]
L'ensemble \(\K[X]\) est un \(\K\)-espace vectoriel.  
La famille :
\[
(1,X,X^2,\dots)
\]
en est une base.
\end{theorembox}

\begin{proof}
Toute combinaison linéaire finie de puissances de \(X\) est un polynôme.  
Réciproquement, tout polynôme s'écrit comme combinaison linéaire finie de :
\[
1,X,X^2,\dots
\]
L'écriture est unique car deux polynômes sont égaux si et seulement si leurs coefficients sont égaux.  
La famille \((1,X,X^2,\dots)\) est donc une base de \(\K[X]\).
\end{proof}

\begin{definitionbox}[Espace \(\K_n[X]\)]
Pour \(n\in\N\), on note :
\[
\K_n[X]=\{P\in\K[X]:\deg(P)\leq n\}.
\]
C'est l'espace vectoriel des polynômes de degré inférieur ou égal à \(n\).
\end{definitionbox}

\begin{theorembox}
L'espace \(\K_n[X]\) est un espace vectoriel de dimension \(n+1\), de base :
\[
(1,X,X^2,\dots,X^n).
\]
\end{theorembox}

\begin{proof}
Tout polynôme de degré inférieur ou égal à \(n\) s'écrit de manière unique :
\[
P(X)=a_0+a_1X+\cdots+a_nX^n.
\]
Donc :
\[
(1,X,\dots,X^n)
\]
est une base de \(\K_n[X]\).  
Elle contient \(n+1\) éléments, donc :
\[
\dim(\K_n[X])=n+1.
\]
\end{proof}

%=========================================================
\section{Degré, coefficient dominant et valuation}

\subsection{Degré}

\begin{definitionbox}[Degré]
Soit \(P\in\K[X]\), \(P\neq0\).  
Si :
\[
P(X)=a_0+a_1X+\cdots+a_nX^n
\]
avec \(a_n\neq0\), alors le degré de \(P\) est :
\[
\deg(P)=n.
\]
\end{definitionbox}

\begin{definitionbox}[Convention]
Par convention :
\[
\deg(0)=-\infty.
\]
\end{definitionbox}

\begin{definitionbox}[Coefficient dominant]
Si \(P\neq0\), le coefficient de \(X^{\deg(P)}\) est appelé le \textbf{coefficient dominant} de \(P\).
\end{definitionbox}

\begin{definitionbox}[Polynôme unitaire]
Un polynôme non nul est dit \textbf{unitaire} si son coefficient dominant vaut \(1\).
\end{definitionbox}

\begin{exercicebox}[Application directe]
Déterminer le degré et le coefficient dominant de :
\[
P(X)=5X^7-3X^2+X-9.
\]
\end{exercicebox}

\begin{correctionbox}
Le terme de plus haut degré est :
\[
5X^7.
\]
Donc :
\[
\deg(P)=7.
\]
Le coefficient dominant est :
\[
5.
\]
Le polynôme n'est pas unitaire.
\end{correctionbox}

\subsection{Propriétés du degré}

\begin{theorembox}[Degré d'une somme]
Pour tous \(P,Q\in\K[X]\),
\[
\deg(P+Q)\leq \max(\deg(P),\deg(Q)).
\]
\end{theorembox}

\begin{proof}
Les coefficients de \(P+Q\) s'obtiennent en additionnant les coefficients de même degré.  
Un terme de degré strictement supérieur à \(\max(\deg(P),\deg(Q))\) ne peut pas apparaître, puisque ni \(P\) ni \(Q\) n'en contient.  
Il peut en revanche y avoir simplification du terme dominant si les degrés sont égaux.  
Donc :
\[
\deg(P+Q)\leq \max(\deg(P),\deg(Q)).
\]
\end{proof}

\begin{erreurbox}
On n'a pas toujours :
\[
\deg(P+Q)=\max(\deg(P),\deg(Q)).
\]
Exemple :
\[
P=X^2+1,\qquad Q=-X^2+X.
\]
Alors :
\[
P+Q=X+1,
\]
donc :
\[
\deg(P+Q)=1<2.
\]
\end{erreurbox}

\begin{theorembox}[Degré d'un produit]
Pour tous \(P,Q\in\K[X]\),
\[
\deg(PQ)=\deg(P)+\deg(Q).
\]
\end{theorembox}

\begin{proof}
Si \(P=0\) ou \(Q=0\), la formule est vraie avec la convention \(\deg(0)=-\infty\).

Supposons \(P\neq0\) et \(Q\neq0\).  
Écrivons :
\[
P(X)=a_nX^n+\text{termes de degré inférieur},
\qquad a_n\neq0,
\]
\[
Q(X)=b_mX^m+\text{termes de degré inférieur},
\qquad b_m\neq0.
\]
Alors le terme de plus haut degré de \(PQ\) est :
\[
a_nb_mX^{n+m}.
\]
Comme \(\K\) est un corps et que \(a_n,b_m\neq0\), on a :
\[
a_nb_m\neq0.
\]
Donc :
\[
\deg(PQ)=n+m=\deg(P)+\deg(Q).
\]
\end{proof}

\begin{corollaire}
Si \(P,Q\in\K[X]\) sont non nuls, alors :
\[
PQ\neq0.
\]
\end{corollaire}

\begin{proof}
Si \(P,Q\neq0\), alors :
\[
\deg(PQ)=\deg(P)+\deg(Q)\in\N.
\]
Donc \(PQ\neq0\).
\end{proof}

\begin{corollaire}[Intégrité de \(\K[X]\)]
L'anneau \(\K[X]\) est intègre : il ne possède pas de diviseurs de zéro.
\end{corollaire}

\begin{retenirbox}
Dans \(\K[X]\), le degré transforme le produit en somme :
\[
\deg(PQ)=\deg(P)+\deg(Q).
\]
C'est l'une des raisons pour lesquelles le degré est un outil fondamental.
\end{retenirbox}

%=========================================================
\section{Polynômes constants et polynômes inversibles}

\begin{definitionbox}[Polynôme constant]
Un polynôme est constant s'il est de degré \(0\) ou s'il est le polynôme nul.
\end{definitionbox}

\begin{theorembox}[Polynômes inversibles de \(\K[X]\)]
Les éléments inversibles de \(\K[X]\) sont exactement les polynômes constants non nuls.
\end{theorembox}

\begin{proof}
Soit \(P\in\K[X]\) inversible.  
Il existe \(Q\in\K[X]\) tel que :
\[
PQ=1.
\]
En prenant les degrés :
\[
\deg(PQ)=\deg(1)=0.
\]
Donc :
\[
\deg(P)+\deg(Q)=0.
\]
Comme \(\deg(P),\deg(Q)\in\N\), on obtient :
\[
\deg(P)=0.
\]
Donc \(P\) est constant non nul.

Réciproquement, si \(P=a\in\K^*\), alors :
\[
P^{-1}=\frac1a.
\]
Donc \(P\) est inversible dans \(\K[X]\).
\end{proof}

\begin{retenirbox}
Dans \(\K[X]\), les seuls inversibles sont les constantes non nulles.
\end{retenirbox}

%=========================================================
\section{Évaluation et fonctions polynomiales}

\subsection{Évaluation}

\begin{definitionbox}[Évaluation en un scalaire]
Soit :
\[
P(X)=a_0+a_1X+\cdots+a_nX^n\in\K[X].
\]
Pour \(\alpha\in\K\), on définit :
\[
P(\alpha)=a_0+a_1\alpha+\cdots+a_n\alpha^n.
\]
\end{definitionbox}

\begin{theorembox}
Pour tout \(\alpha\in\K\), l'application :
\[
\operatorname{ev}_\alpha:\K[X]\to\K,
\qquad
P\mapsto P(\alpha),
\]
est un morphisme d'anneaux.
\end{theorembox}

\begin{proof}
Soient \(P,Q\in\K[X]\).  
Par définition de l'addition :
\[
(P+Q)(\alpha)=P(\alpha)+Q(\alpha).
\]
Par définition du produit :
\[
(PQ)(\alpha)=P(\alpha)Q(\alpha).
\]
Enfin :
\[
1(\alpha)=1.
\]
Donc \(\operatorname{ev}_\alpha\) est un morphisme d'anneaux.
\end{proof}

\subsection{Polynôme et fonction polynomiale}

À tout polynôme \(P\in\K[X]\), on associe une fonction :
\[
\widetilde P:\K\to\K,
\qquad
\alpha\mapsto P(\alpha).
\]

\begin{theorembox}
Si \(\K\) est infini, deux polynômes de \(\K[X]\) définissant la même fonction polynomiale sont égaux.
\end{theorembox}

\begin{proof}
Soient \(P,Q\in\K[X]\) tels que :
\[
\forall x\in\K,\qquad P(x)=Q(x).
\]
Alors :
\[
R=P-Q
\]
s'annule en tout élément de \(\K\).  
Si \(R\neq0\), il aurait un degré \(n\), donc au plus \(n\) racines distinctes.  
Mais \(\K\) est infini, donc \(R\) aurait une infinité de racines.  
Contradiction.  
Donc :
\[
R=0,
\]
et :
\[
P=Q.
\]
\end{proof}

\begin{erreurbox}
Sur un corps fini, deux polynômes différents peuvent définir la même fonction.  
Par exemple sur \(\mathbb F_2\), les polynômes \(X^2\) et \(X\) donnent la même fonction, car :
\[
0^2=0,\qquad 1^2=1.
\]
Mais comme polynômes :
\[
X^2\neq X.
\]
\end{erreurbox}

%=========================================================
\section{Division euclidienne dans \(\K[X]\)}

\subsection{Énoncé}

\begin{theorembox}[Division euclidienne]
Soient \(A,B\in\K[X]\), avec \(B\neq0\).  
Il existe un unique couple \((Q,R)\in\K[X]^2\) tel que :
\[
A=BQ+R
\]
et :
\[
\deg(R)<\deg(B).
\]
Le polynôme \(Q\) est le quotient, et \(R\) est le reste.
\end{theorembox}

\subsection{Unicité}

\begin{proof}
Supposons qu'il existe deux écritures :
\[
A=BQ+R,
\qquad
A=BQ'+R',
\]
avec :
\[
\deg(R)<\deg(B),
\qquad
\deg(R')<\deg(B).
\]
En soustrayant :
\[
B(Q-Q')=R'-R.
\]
Si \(Q-Q'\neq0\), alors :
\[
\deg(B(Q-Q'))=\deg(B)+\deg(Q-Q')\geq\deg(B).
\]
Mais :
\[
\deg(R'-R)\leq\max(\deg(R'),\deg(R))<\deg(B).
\]
Contradiction.  
Donc :
\[
Q-Q'=0.
\]
Ainsi :
\[
Q=Q',
\]
puis :
\[
R=R'.
\]
\end{proof}

\subsection{Existence}

\begin{proof}
On raisonne par récurrence sur \(\deg(A)\).

Si \(A=0\), on prend :
\[
Q=0,\qquad R=0.
\]

Supposons \(A\neq0\).  
Si :
\[
\deg(A)<\deg(B),
\]
on prend :
\[
Q=0,\qquad R=A.
\]

Sinon, posons :
\[
\deg(A)=n,
\qquad
\deg(B)=m,
\qquad
n\geq m.
\]
Écrivons les termes dominants :
\[
A=a_nX^n+\cdots,
\qquad
B=b_mX^m+\cdots,
\]
avec \(a_n,b_m\neq0\).

Considérons :
\[
A_1=A-\frac{a_n}{b_m}X^{n-m}B.
\]
Le terme de degré \(n\) disparaît dans \(A_1\).  
Donc :
\[
\deg(A_1)<\deg(A).
\]
Par hypothèse de récurrence, il existe \(Q_1,R\) tels que :
\[
A_1=BQ_1+R,
\qquad
\deg(R)<\deg(B).
\]
Ainsi :
\[
A
=
A_1+\frac{a_n}{b_m}X^{n-m}B
=
B\left(Q_1+\frac{a_n}{b_m}X^{n-m}\right)+R.
\]
On pose :
\[
Q=Q_1+\frac{a_n}{b_m}X^{n-m}.
\]
On a bien :
\[
A=BQ+R,
\qquad
\deg(R)<\deg(B).
\]
\end{proof}

\begin{methodbox}[Effectuer une division euclidienne]
Pour diviser \(A\) par \(B\neq0\) :
\begin{enumerate}
    \item comparer les termes dominants ;
    \item choisir un monôme qui annule le terme dominant de \(A\) ;
    \item soustraire le multiple correspondant de \(B\) ;
    \item recommencer jusqu'à obtenir un reste de degré strictement inférieur à \(\deg(B)\).
\end{enumerate}
\end{methodbox}

\begin{exercicebox}[Division euclidienne]
Effectuer la division euclidienne de :
\[
A(X)=X^3-2X+1
\]
par :
\[
B(X)=X-1.
\]
\end{exercicebox}

\begin{correctionbox}
On cherche \(Q,R\) tels que :
\[
X^3-2X+1=(X-1)Q+R,
\qquad
\deg(R)<1.
\]
Le reste est donc constant.

On divise :
\[
X^3-2X+1=(X-1)(X^2+X-1)+0.
\]
En effet :
\[
(X-1)(X^2+X-1)=X^3+X^2-X-X^2-X+1=X^3-2X+1.
\]
Donc :
\[
Q=X^2+X-1,
\qquad
R=0.
\]
\end{correctionbox}

%=========================================================
\section{Racines d'un polynôme}

\subsection{Définition}

\begin{definitionbox}[Racine]
Soit \(P\in\K[X]\).  
Un scalaire \(\alpha\in\K\) est une racine de \(P\) si :
\[
P(\alpha)=0.
\]
\end{definitionbox}

\subsection{Théorème des racines}

\begin{theorembox}[Théorème des racines]
Soit \(P\in\K[X]\) et \(\alpha\in\K\).  
Alors :
\[
P(\alpha)=0
\Longleftrightarrow
X-\alpha \text{ divise } P.
\]
\end{theorembox}

\begin{proof}
On effectue la division euclidienne de \(P\) par \(X-\alpha\).  
Il existe \(Q\in\K[X]\) et \(r\in\K\) tels que :
\[
P(X)=(X-\alpha)Q(X)+r.
\]
En évaluant en \(\alpha\), on obtient :
\[
P(\alpha)=(\alpha-\alpha)Q(\alpha)+r=r.
\]
Donc :
\[
r=P(\alpha).
\]
Ainsi :
\[
P(\alpha)=0
\]
si et seulement si :
\[
r=0,
\]
si et seulement si :
\[
P(X)=(X-\alpha)Q(X).
\]
Donc \(X-\alpha\) divise \(P\).
\end{proof}

\begin{retenirbox}
Chercher une racine revient à chercher un facteur de degré \(1\).
\[
\alpha\text{ racine de }P
\Longleftrightarrow
X-\alpha\mid P.
\]
\end{retenirbox}

\subsection{Nombre maximal de racines}

\begin{theorembox}
Un polynôme non nul de degré \(n\) possède au plus \(n\) racines distinctes dans \(\K\).
\end{theorembox}

\begin{proof}
On raisonne par récurrence sur \(n\).

Pour \(n=0\), un polynôme constant non nul n'a aucune racine.

Supposons le résultat vrai pour les polynômes de degré \(n-1\).  
Soit \(P\) un polynôme de degré \(n\).

Si \(P\) n'a pas de racine, le résultat est évident.  
Sinon, soit \(\alpha\) une racine de \(P\).  
Alors :
\[
P=(X-\alpha)Q,
\]
avec :
\[
\deg(Q)=n-1.
\]
Toute racine de \(P\) différente de \(\alpha\) est une racine de \(Q\).  
En effet, si \(\beta\neq\alpha\) et \(P(\beta)=0\), alors :
\[
0=(\beta-\alpha)Q(\beta).
\]
Comme \(\beta-\alpha\neq0\), on a :
\[
Q(\beta)=0.
\]
Par hypothèse de récurrence, \(Q\) possède au plus \(n-1\) racines distinctes.  
Donc \(P\) possède au plus \(n\) racines distinctes.
\end{proof}

\begin{corollaire}
Si un polynôme \(P\) de degré inférieur ou égal à \(n\) possède \(n+1\) racines distinctes, alors :
\[
P=0.
\]
\end{corollaire}

\begin{corollaire}
Si deux polynômes de degré inférieur ou égal à \(n\) coïncident en \(n+1\) points distincts, alors ils sont égaux.
\end{corollaire}

\begin{proof}
Soient \(P,Q\in\K_n[X]\) tels que :
\[
P(a_i)=Q(a_i)
\]
pour \(n+1\) scalaires distincts \(a_0,\dots,a_n\).  
Alors \(P-Q\) possède \(n+1\) racines distinctes et :
\[
\deg(P-Q)\leq n.
\]
Donc :
\[
P-Q=0.
\]
Ainsi :
\[
P=Q.
\]
\end{proof}

\begin{erreurbox}
Un polynôme de degré \(n\) possède au plus \(n\) racines dans \(\K\), mais il n'en possède pas forcément \(n\).  
Exemple :
\[
X^2+1
\]
n'a aucune racine réelle, mais possède deux racines complexes.
\end{erreurbox}

%=========================================================
\section{Multiplicité d'une racine}

\subsection{Définition}

\begin{definitionbox}[Racine de multiplicité \(m\)]
Soit \(P\in\K[X]\), \(P\neq0\), et soit \(\alpha\in\K\).  
On dit que \(\alpha\) est racine de multiplicité \(m\geq1\) de \(P\) si :
\[
P(X)=(X-\alpha)^mQ(X),
\]
avec :
\[
Q(\alpha)\neq0.
\]
\end{definitionbox}

Autrement dit, \(m\) est le plus grand entier tel que :
\[
(X-\alpha)^m\mid P.
\]

\begin{definitionbox}
Une racine de multiplicité \(1\) est dite \textbf{simple}.  
Une racine de multiplicité au moins \(2\) est dite \textbf{multiple}.
\end{definitionbox}

\begin{exercicebox}[Multiplicités]
Déterminer les racines et leurs multiplicités pour :
\[
P(X)=(X-1)^3(X+2)^2(X^2+1).
\]
dans \(\R[X]\), puis dans \(\C[X]\).
\end{exercicebox}

\begin{correctionbox}
Dans \(\R[X]\), les racines réelles sont :
\[
1 \quad\text{de multiplicité }3,
\]
et :
\[
-2 \quad\text{de multiplicité }2.
\]
Le facteur \(X^2+1\) n'a pas de racine réelle.

Dans \(\C[X]\), on a :
\[
X^2+1=(X-i)(X+i).
\]
Donc les racines complexes sont :
\[
1 \text{ de multiplicité }3,
\qquad
-2 \text{ de multiplicité }2,
\qquad
i \text{ de multiplicité }1,
\qquad
-i \text{ de multiplicité }1.
\]
\end{correctionbox}

\subsection{Somme des multiplicités}

\begin{theorembox}
La somme des multiplicités des racines distinctes d'un polynôme non nul est inférieure ou égale à son degré.
\end{theorembox}

\begin{proof}
Soient \(\alpha_1,\dots,\alpha_r\) les racines distinctes de \(P\), de multiplicités respectives :
\[
m_1,\dots,m_r.
\]
Alors :
\[
(X-\alpha_1)^{m_1}\cdots(X-\alpha_r)^{m_r}
\]
divise \(P\).  
Donc il existe \(Q\in\K[X]\) tel que :
\[
P(X)=\prod_{i=1}^r(X-\alpha_i)^{m_i}Q(X).
\]
En prenant les degrés :
\[
\deg(P)=m_1+\cdots+m_r+\deg(Q).
\]
Comme \(\deg(Q)\geq0\) si \(Q\neq0\), on obtient :
\[
m_1+\cdots+m_r\leq\deg(P).
\]
\end{proof}

%=========================================================
\section{Dérivée formelle et formule de Taylor}

\subsection{Dérivée formelle}

\begin{definitionbox}[Dérivée formelle]
Soit :
\[
P(X)=a_0+a_1X+\cdots+a_nX^n.
\]
La dérivée formelle de \(P\) est :
\[
P'(X)=a_1+2a_2X+\cdots+na_nX^{n-1}.
\]
\end{definitionbox}

\begin{theorembox}[Propriétés]
Pour tous \(P,Q\in\K[X]\) et tout \(\lambda\in\K\),
\[
(P+Q)'=P'+Q',
\]
\[
(\lambda P)'=\lambda P',
\]
\[
(PQ)'=P'Q+PQ'.
\]
\end{theorembox}

\begin{proof}
Les deux premières propriétés se vérifient directement coefficient par coefficient.

Pour le produit, il suffit d'abord de vérifier la formule sur deux monômes :
\[
(X^aX^b)'=(X^{a+b})'=(a+b)X^{a+b-1}.
\]
D'autre part :
\[
(X^a)'X^b+X^a(X^b)'
=
aX^{a-1}X^b+bX^aX^{b-1}
=
(a+b)X^{a+b-1}.
\]
Par bilinéarité du produit et de la dérivation formelle, la formule s'étend à tous les polynômes.
\end{proof}

\subsection{Dérivées successives}

On définit :
\[
P^{(0)}=P,
\]
et, pour \(k\geq1\),
\[
P^{(k)}=(P^{(k-1)})'.
\]

\begin{theorembox}
Si :
\[
P(X)=\sum_{j=0}^n a_jX^j,
\]
alors :
\[
P^{(k)}(X)=\sum_{j=k}^n j(j-1)\cdots(j-k+1)a_jX^{j-k}.
\]
\end{theorembox}

\subsection{Formule de Taylor polynomiale}

\begin{theorembox}[Formule de Taylor pour les polynômes]
Soit \(P\in\K[X]\), où \(\K\) est de caractéristique nulle, et soit \(a\in\K\).  
Alors :
\[
P(X)=\sum_{k=0}^{\deg(P)}
\frac{P^{(k)}(a)}{k!}(X-a)^k.
\]
\end{theorembox}

\begin{proof}
Les polynômes :
\[
1,\ (X-a),\ (X-a)^2,\dots,(X-a)^n
\]
forment une base de \(\K_n[X]\).  
Donc il existe des scalaires \(c_0,\dots,c_n\) tels que :
\[
P(X)=\sum_{k=0}^n c_k(X-a)^k.
\]
Évaluons en \(X=a\).  
On obtient :
\[
P(a)=c_0.
\]

Dérivons :
\[
P'(X)=\sum_{k=1}^n kc_k(X-a)^{k-1}.
\]
En évaluant en \(a\), on obtient :
\[
P'(a)=c_1.
\]

Plus généralement, la dérivée \(j\)-ième donne :
\[
P^{(j)}(a)=j!c_j.
\]
Donc :
\[
c_j=\frac{P^{(j)}(a)}{j!}.
\]
Ainsi :
\[
P(X)=\sum_{k=0}^n
\frac{P^{(k)}(a)}{k!}(X-a)^k.
\]
\end{proof}

\subsection{Lien avec les racines multiples}

\begin{theorembox}[Critère différentiel de multiplicité]
Soit \(P\in\K[X]\), où \(\K\) est de caractéristique nulle.  
Le scalaire \(\alpha\) est racine de multiplicité au moins \(m\) de \(P\) si et seulement si :
\[
P(\alpha)=P'(\alpha)=\cdots=P^{(m-1)}(\alpha)=0.
\]
Il est racine de multiplicité exactement \(m\) si, de plus :
\[
P^{(m)}(\alpha)\neq0.
\]
\end{theorembox}

\begin{proof}
Par la formule de Taylor en \(\alpha\),
\[
P(X)=\sum_{k=0}^{n}
\frac{P^{(k)}(\alpha)}{k!}(X-\alpha)^k.
\]
Le polynôme est divisible par \((X-\alpha)^m\) si et seulement si les coefficients des puissances :
\[
(X-\alpha)^0,\dots,(X-\alpha)^{m-1}
\]
sont nuls.  
Cela équivaut à :
\[
P(\alpha)=P'(\alpha)=\cdots=P^{(m-1)}(\alpha)=0.
\]
La multiplicité est exactement \(m\) si le premier coefficient non nul est celui de \((X-\alpha)^m\), donc si :
\[
P^{(m)}(\alpha)\neq0.
\]
\end{proof}

\begin{methodbox}[Détecter une racine multiple]
Pour montrer que \(\alpha\) est racine multiple de \(P\) :
\begin{enumerate}
    \item vérifier que \(P(\alpha)=0\) ;
    \item vérifier que \(P'(\alpha)=0\).
\end{enumerate}
Pour connaître la multiplicité exacte, on continue avec les dérivées successives.
\end{methodbox}

%=========================================================
\section{Factorisation sur \(\C\) et sur \(\R\)}

\subsection{Polynômes scindés}

\begin{definitionbox}[Polynôme scindé]
Un polynôme \(P\in\K[X]\) est scindé sur \(\K\) s'il peut s'écrire comme produit de facteurs de degré \(1\) :
\[
P(X)=a\prod_{i=1}^r(X-\alpha_i)^{m_i},
\]
avec :
\[
a\in\K^*,
\qquad
\alpha_i\in\K.
\]
\end{definitionbox}

\subsection{Factorisation sur \(\C\)}

\begin{theorembox}[Théorème fondamental de l'algèbre]
Tout polynôme non constant de \(\C[X]\) possède au moins une racine complexe.
\end{theorembox}

\begin{corollaire}
Tout polynôme non constant de \(\C[X]\) est scindé sur \(\C\).
\end{corollaire}

\begin{proof}
Soit \(P\in\C[X]\), non constant.  
Par le théorème fondamental de l'algèbre, \(P\) possède une racine \(\alpha_1\).  
Donc :
\[
P=(X-\alpha_1)Q_1.
\]
Si \(Q_1\) est non constant, il possède une racine \(\alpha_2\), et on recommence.  
Comme le degré diminue à chaque étape, le processus s'arrête.  
On obtient une factorisation complète en facteurs de degré \(1\).
\end{proof}

\subsection{Factorisation sur \(\R\)}

\begin{theorembox}
Tout polynôme non constant de \(\R[X]\) se factorise comme produit de polynômes irréductibles de degré \(1\) ou de degré \(2\).
\end{theorembox}

\begin{proof}
Soit \(P\in\R[X]\).  
On le considère comme un polynôme de \(\C[X]\).  
Sur \(\C\), il est scindé :
\[
P(X)=a\prod_{k=1}^n(X-z_k).
\]
Comme les coefficients de \(P\) sont réels, les racines complexes non réelles apparaissent par paires conjuguées.  
Si \(z\notin\R\) est racine, alors \(\overline z\) est aussi racine.

Le produit associé est :
\[
(X-z)(X-\overline z)
=
X^2-(z+\overline z)X+z\overline z.
\]
Or :
\[
z+\overline z=2\operatorname{Re}(z)\in\R,
\]
et :
\[
z\overline z=|z|^2\in\R.
\]
Ce facteur est donc un polynôme réel de degré \(2\).  
Les racines réelles donnent des facteurs de degré \(1\).  
On obtient donc une factorisation dans \(\R[X]\) en facteurs de degré \(1\) et \(2\).
\end{proof}

\begin{exercicebox}[Factorisation de \(X^4+1\)]
Factoriser :
\[
X^4+1
\]
dans \(\C[X]\), puis dans \(\R[X]\).
\end{exercicebox}

\begin{correctionbox}
Dans \(\C[X]\), on cherche les racines de :
\[
X^4=-1=e^{i\pi}.
\]
Les racines sont :
\[
e^{i\pi/4},
\quad
e^{3i\pi/4},
\quad
e^{5i\pi/4},
\quad
e^{7i\pi/4}.
\]
Donc \(X^4+1\) est scindé sur \(\C\).

Dans \(\R[X]\), on regroupe les racines conjuguées.  
On obtient :
\[
X^4+1=(X^2+\sqrt2X+1)(X^2-\sqrt2X+1).
\]
\end{correctionbox}

%=========================================================
\section{Polynômes irréductibles}

\begin{definitionbox}[Polynôme irréductible]
Un polynôme \(P\in\K[X]\) est irréductible sur \(\K\) si :
\begin{enumerate}
    \item \(P\) est non constant ;
    \item si \(P=AB\), alors \(A\) ou \(B\) est constant.
\end{enumerate}
\end{definitionbox}

\begin{retenirbox}
Un polynôme irréductible est un polynôme que l'on ne peut pas factoriser non trivialement dans \(\K[X]\).
\end{retenirbox}

\begin{theorembox}
Dans \(\C[X]\), les polynômes irréductibles sont exactement les polynômes de degré \(1\).
\end{theorembox}

\begin{proof}
Tout polynôme non constant de degré au moins \(2\) possède une racine complexe, donc est divisible par un facteur de degré \(1\).  
Il est donc réductible.  
Les polynômes de degré \(1\), eux, sont irréductibles.
\end{proof}

\begin{theorembox}
Dans \(\R[X]\), les polynômes irréductibles sont :
\begin{itemize}
    \item les polynômes de degré \(1\) ;
    \item les polynômes de degré \(2\) sans racine réelle.
\end{itemize}
\end{theorembox}

\begin{proof}
Un polynôme de degré \(1\) est irréductible.

Un polynôme réel de degré \(2\) sans racine réelle est irréductible.  
En effet, s'il était réductible, il serait produit de deux polynômes non constants, donc de deux polynômes de degré \(1\), ce qui donnerait une racine réelle.

Réciproquement, soit \(P\in\R[X]\) irréductible.  
D'après la factorisation réelle, \(P\) se factorise en produits de facteurs de degré \(1\) et de facteurs quadratiques irréductibles.  
Comme \(P\) est irréductible, il ne peut contenir qu'un seul facteur.  
Donc son degré est \(1\) ou \(2\).  
Dans le cas du degré \(2\), il ne doit pas avoir de racine réelle.
\end{proof}

\begin{theorembox}
Un polynôme réel de degré \(2\) ou \(3\) est irréductible sur \(\R\) si et seulement s'il n'a pas de racine réelle.
\end{theorembox}

\begin{proof}
Si \(P\) a une racine réelle \(\alpha\), alors :
\[
X-\alpha
\]
divise \(P\).  
Donc \(P\) est réductible.

Réciproquement, si \(P\) est réductible et de degré \(2\) ou \(3\), alors :
\[
P=AB,
\]
avec \(A,B\) non constants.  
La somme des degrés vaut \(2\) ou \(3\).  
L'un des facteurs est donc de degré \(1\).  
Un facteur de degré \(1\) donne une racine réelle.  
Donc si \(P\) n'a pas de racine réelle, il est irréductible.
\end{proof}

%=========================================================
\section{Arithmétique dans \(\K[X]\)}

\subsection{Divisibilité}

\begin{definitionbox}[Divisibilité]
Soient \(A,B\in\K[X]\).  
On dit que \(B\) divise \(A\), et on note :
\[
B\mid A,
\]
s'il existe \(Q\in\K[X]\) tel que :
\[
A=BQ.
\]
\end{definitionbox}

\subsection{PGCD}

\begin{definitionbox}[PGCD]
Soient \(A,B\in\K[X]\), non tous deux nuls.  
Un PGCD de \(A\) et \(B\) est un polynôme \(D\) tel que :
\begin{enumerate}
    \item \(D\mid A\) et \(D\mid B\) ;
    \item tout diviseur commun de \(A\) et \(B\) divise \(D\).
\end{enumerate}
\end{definitionbox}

Le PGCD est défini à multiplication par une constante non nulle près.  
On choisit généralement le PGCD unitaire.

\subsection{Algorithme d'Euclide}

\begin{methodbox}[Algorithme d'Euclide]
Pour calculer \(\pgcd(A,B)\), on effectue des divisions euclidiennes successives :
\[
A=BQ_1+R_1,
\]
\[
B=R_1Q_2+R_2,
\]
\[
R_1=R_2Q_3+R_3,
\]
et ainsi de suite, jusqu'à obtenir un reste nul.  
Le dernier reste non nul est un PGCD de \(A\) et \(B\).
\end{methodbox}

\begin{theorembox}
L'algorithme d'Euclide fournit un PGCD de deux polynômes non tous deux nuls.
\end{theorembox}

\begin{proof}
À chaque division euclidienne :
\[
A=BQ+R,
\]
les diviseurs communs de \(A\) et \(B\) sont exactement les diviseurs communs de \(B\) et \(R\).  
En effet :
\[
R=A-BQ.
\]
Donc tout diviseur commun de \(A\) et \(B\) divise \(R\).  
Réciproquement :
\[
A=BQ+R.
\]
Donc tout diviseur commun de \(B\) et \(R\) divise \(A\).  
Ainsi le PGCD ne change pas à chaque étape.

Comme les degrés des restes diminuent strictement, l'algorithme s'arrête.  
Le dernier reste non nul est alors un PGCD.
\end{proof}

\subsection{Bézout}

\begin{theorembox}[Théorème de Bézout]
Soient \(A,B\in\K[X]\), non tous deux nuls.  
Si \(D\) est un PGCD de \(A\) et \(B\), alors il existe \(U,V\in\K[X]\) tels que :
\[
AU+BV=D.
\]
En particulier, \(A\) et \(B\) sont premiers entre eux si et seulement s'il existe \(U,V\in\K[X]\) tels que :
\[
AU+BV=1.
\]
\end{theorembox}

\begin{proof}
L'algorithme d'Euclide exprime chaque reste comme combinaison linéaire des deux précédents.  
En remontant les égalités, on exprime le dernier reste non nul \(D\) sous la forme :
\[
D=AU+BV.
\]
Si \(A\) et \(B\) sont premiers entre eux, on choisit \(D=1\), donc :
\[
AU+BV=1.
\]

Réciproquement, s'il existe \(U,V\) tels que :
\[
AU+BV=1,
\]
tout diviseur commun de \(A\) et \(B\) divise le membre de gauche, donc divise \(1\).  
Ainsi les seuls diviseurs communs sont les constantes non nulles.  
Donc \(A\) et \(B\) sont premiers entre eux.
\end{proof}

\begin{retenirbox}
Bézout transforme une propriété de divisibilité en une identité :
\[
A\wedge B=1
\Longleftrightarrow
\exists U,V,\quad AU+BV=1.
\]
\end{retenirbox}

\subsection{Lemme de Gauss}

\begin{theorembox}[Lemme de Gauss]
Si \(A,B,C\in\K[X]\), si \(A\mid BC\), et si \(A\) est premier avec \(B\), alors :
\[
A\mid C.
\]
\end{theorembox}

\begin{proof}
Comme \(A\) et \(B\) sont premiers entre eux, il existe \(U,V\in\K[X]\) tels que :
\[
AU+BV=1.
\]
En multipliant par \(C\), on obtient :
\[
ACU+BCV=C.
\]
Or \(A\mid ACU\), et par hypothèse \(A\mid BC\), donc \(A\mid BCV\).  
Ainsi \(A\) divise la somme :
\[
ACU+BCV=C.
\]
Donc :
\[
A\mid C.
\]
\end{proof}

%=========================================================
\section{Interpolation de Lagrange}

\subsection{Problème d'interpolation}

On se donne \(n+1\) points d'abscisses distinctes :
\[
a_0,\dots,a_n\in\K,
\]
et \(n+1\) valeurs :
\[
y_0,\dots,y_n\in\K.
\]
On cherche un polynôme \(P\in\K_n[X]\) tel que :
\[
P(a_i)=y_i
\]
pour tout \(i\in\{0,\dots,n\}\).

\begin{theorembox}[Existence et unicité]
Soient \(a_0,\dots,a_n\in\K\) deux à deux distincts.  
Pour tous \(y_0,\dots,y_n\in\K\), il existe un unique polynôme \(P\in\K_n[X]\) tel que :
\[
P(a_i)=y_i
\]
pour tout \(i\).
\end{theorembox}

\begin{proof}
Considérons l'application :
\[
\Phi:\K_n[X]\to\K^{n+1},
\qquad
P\mapsto(P(a_0),\dots,P(a_n)).
\]
Cette application est linéaire.

Montrons qu'elle est injective.  
Si :
\[
\Phi(P)=0,
\]
alors :
\[
P(a_0)=\cdots=P(a_n)=0.
\]
Donc \(P\) possède \(n+1\) racines distinctes.  
Or :
\[
\deg(P)\leq n.
\]
Donc :
\[
P=0.
\]
Ainsi \(\Phi\) est injective.

Comme :
\[
\dim(\K_n[X])=n+1
\]
et :
\[
\dim(\K^{n+1})=n+1,
\]
l'application \(\Phi\) est un isomorphisme.  
Elle est donc bijective, ce qui donne existence et unicité.
\end{proof}

\subsection{Polynômes de Lagrange}

\begin{definitionbox}[Polynômes interpolateurs de Lagrange]
Pour \(i\in\{0,\dots,n\}\), on définit :
\[
L_i(X)=\prod_{\substack{0\leq j\leq n\\ j\neq i}}
\frac{X-a_j}{a_i-a_j}.
\]
\end{definitionbox}

\begin{theorembox}
Les polynômes \(L_i\) vérifient :
\[
L_i(a_j)=\delta_{ij}.
\]
\end{theorembox}

\begin{proof}
Si \(j=i\), alors :
\[
L_i(a_i)=
\prod_{k\neq i}
\frac{a_i-a_k}{a_i-a_k}=1.
\]
Si \(j\neq i\), alors dans le produit définissant \(L_i(a_j)\), il y a un facteur :
\[
a_j-a_j=0.
\]
Donc :
\[
L_i(a_j)=0.
\]
Ainsi :
\[
L_i(a_j)=\delta_{ij}.
\]
\end{proof}

\begin{theorembox}[Formule d'interpolation de Lagrange]
L'unique polynôme \(P\in\K_n[X]\) tel que :
\[
P(a_i)=y_i
\]
pour tout \(i\) est :
\[
P(X)=\sum_{i=0}^n y_iL_i(X).
\]
\end{theorembox}

\begin{proof}
Posons :
\[
P(X)=\sum_{i=0}^n y_iL_i(X).
\]
Pour \(j\in\{0,\dots,n\}\),
\[
P(a_j)=\sum_{i=0}^n y_iL_i(a_j).
\]
Or :
\[
L_i(a_j)=\delta_{ij}.
\]
Donc :
\[
P(a_j)=y_j.
\]
Par unicité du polynôme interpolateur, c'est le seul.
\end{proof}

\begin{retenirbox}
Les polynômes \(L_i\) forment une base de \(\K_n[X]\).  
Dans cette base :
\[
P=\sum_{i=0}^n P(a_i)L_i.
\]
\end{retenirbox}

%=========================================================
\section{Déterminant de Vandermonde}

\begin{definitionbox}[Matrice de Vandermonde]
Soient \(a_0,\dots,a_n\in\K\).  
La matrice de Vandermonde associée est :
\[
V=
\begin{pmatrix}
1&a_0&a_0^2&\cdots&a_0^n\\
1&a_1&a_1^2&\cdots&a_1^n\\
\vdots&\vdots&\vdots&&\vdots\\
1&a_n&a_n^2&\cdots&a_n^n
\end{pmatrix}.
\]
\end{definitionbox}

\begin{theorembox}[Déterminant de Vandermonde]
On a :
\[
\det(V)=\prod_{0\leq i<j\leq n}(a_j-a_i).
\]
\end{theorembox}

\begin{proof}
On considère \(\det(V)\) comme un polynôme en les variables \(a_0,\dots,a_n\).

Si deux \(a_i\) sont égaux, alors deux lignes de \(V\) sont égales.  
Donc :
\[
\det(V)=0.
\]
Ainsi, pour tout \(i<j\), le facteur \(a_j-a_i\) divise \(\det(V)\).

Le produit :
\[
\prod_{0\leq i<j\leq n}(a_j-a_i)
\]
a le même degré total que \(\det(V)\), à savoir :
\[
0+1+\cdots+n=\frac{n(n+1)}2.
\]
Donc :
\[
\det(V)=C\prod_{0\leq i<j\leq n}(a_j-a_i)
\]
pour une constante \(C\).

Pour déterminer \(C\), on regarde le coefficient du monôme :
\[
a_0^0a_1^1a_2^2\cdots a_n^n.
\]
Dans le déterminant, ce coefficient vaut \(1\), obtenu par la diagonale principale.  
Dans le produit, ce coefficient vaut aussi \(1\).  
Donc :
\[
C=1.
\]
Ainsi :
\[
\det(V)=\prod_{0\leq i<j\leq n}(a_j-a_i).
\]
\end{proof}

\begin{corollaire}
La matrice de Vandermonde est inversible si et seulement si les \(a_i\) sont deux à deux distincts.
\end{corollaire}

%=========================================================
\section{Fractions rationnelles}

\subsection{Définition}

\begin{definitionbox}[Fraction rationnelle]
Une fraction rationnelle sur \(\K\) est une expression de la forme :
\[
\frac{P(X)}{Q(X)},
\]
où :
\[
P,Q\in\K[X],
\qquad
Q\neq0.
\]
\end{definitionbox}

Deux fractions :
\[
\frac{P}{Q}
\quad\text{et}\quad
\frac{P'}{Q'}
\]
sont égales si :
\[
PQ'=P'Q.
\]

\begin{definitionbox}[Partie entière]
Soit \(F=\dfrac{P}{Q}\) une fraction rationnelle.  
Par division euclidienne :
\[
P=QE+R,
\qquad
\deg(R)<\deg(Q).
\]
Alors :
\[
F=E+\frac{R}{Q}.
\]
Le polynôme \(E\) est appelé la partie entière de \(F\).
\end{definitionbox}

\subsection{Exemple de décomposition}

\begin{exercicebox}[Décomposition simple]
Décomposer :
\[
\frac{1}{X(X-1)}
\]
sous la forme :
\[
\frac{a}{X}+\frac{b}{X-1}.
\]
\end{exercicebox}

\begin{correctionbox}
On cherche \(a,b\in\K\) tels que :
\[
\frac{1}{X(X-1)}
=
\frac{a}{X}+\frac{b}{X-1}.
\]
En mettant au même dénominateur :
\[
\frac{1}{X(X-1)}
=
\frac{a(X-1)+bX}{X(X-1)}.
\]
Donc :
\[
1=a(X-1)+bX.
\]
En développant :
\[
1=(a+b)X-a.
\]
On identifie les coefficients :
\[
a+b=0,
\qquad
-a=1.
\]
Donc :
\[
a=-1,
\qquad
b=1.
\]
Ainsi :
\[
\frac{1}{X(X-1)}
=
-\frac1X+\frac1{X-1}.
\]
\end{correctionbox}

%=========================================================
\section{Polynômes et algèbre linéaire}

\subsection{Polynôme d'un endomorphisme}

\begin{definitionbox}[Polynôme d'un endomorphisme]
Soit \(E\) un \(\K\)-espace vectoriel et soit \(u\in\mathcal L(E)\).  
Si :
\[
P(X)=a_0+a_1X+\cdots+a_nX^n,
\]
on définit :
\[
P(u)=a_0\Id_E+a_1u+\cdots+a_nu^n.
\]
\end{definitionbox}

\begin{definitionbox}[Polynôme d'une matrice]
Si \(A\in\M_n(\K)\), on définit :
\[
P(A)=a_0I_n+a_1A+\cdots+a_nA^n.
\]
\end{definitionbox}

\begin{definitionbox}[Polynôme annulateur]
Un polynôme \(P\in\K[X]\) est annulateur de \(u\) si :
\[
P(u)=0.
\]
\end{definitionbox}

\begin{theorembox}
Si \(x\neq0\) est vecteur propre de \(u\) associé à la valeur propre \(\lambda\), alors :
\[
P(u)(x)=P(\lambda)x.
\]
\end{theorembox}

\begin{proof}
Comme :
\[
u(x)=\lambda x,
\]
on montre par récurrence que :
\[
u^k(x)=\lambda^kx.
\]
Donc :
\[
P(u)(x)
=
a_0x+a_1u(x)+\cdots+a_nu^n(x)
=
(a_0+a_1\lambda+\cdots+a_n\lambda^n)x
=
P(\lambda)x.
\]
\end{proof}

\begin{corollaire}
Si \(P(u)=0\), alors toute valeur propre \(\lambda\) de \(u\) vérifie :
\[
P(\lambda)=0.
\]
\end{corollaire}

\begin{proof}
Soit \(x\neq0\) un vecteur propre associé à \(\lambda\).  
Alors :
\[
0=P(u)(x)=P(\lambda)x.
\]
Comme \(x\neq0\), on obtient :
\[
P(\lambda)=0.
\]
\end{proof}

\begin{erreurbox}
Les valeurs propres sont racines de tout polynôme annulateur.  
Mais toutes les racines d'un polynôme annulateur ne sont pas forcément des valeurs propres.
\end{erreurbox}

%=========================================================
\section{Méthodes essentielles}

\begin{methodbox}[Montrer que deux polynômes sont égaux]
Méthodes possibles :
\begin{enumerate}
    \item identifier tous les coefficients ;
    \item montrer que leur différence est nulle ;
    \item si le degré est majoré par \(n\), montrer qu'ils coïncident en \(n+1\) points distincts ;
    \item utiliser une base adaptée, par exemple la base de Lagrange.
\end{enumerate}
\end{methodbox}

\begin{methodbox}[Factoriser un polynôme]
Pour factoriser \(P\) :
\begin{enumerate}
    \item chercher des racines évidentes ;
    \item utiliser le théorème :
    \[
    P(\alpha)=0\Longleftrightarrow X-\alpha\mid P ;
    \]
    \item effectuer une division euclidienne ;
    \item recommencer sur le quotient ;
    \item selon le corps de base, traiter les facteurs irréductibles restants.
\end{enumerate}
\end{methodbox}

\begin{methodbox}[Trouver une multiplicité]
Pour déterminer la multiplicité de \(\alpha\) :
\begin{enumerate}
    \item vérifier que \(P(\alpha)=0\) ;
    \item calculer \(P'(\alpha)\), \(P''(\alpha)\), etc. ;
    \item la multiplicité est le plus petit rang \(m\) tel que :
    \[
    P^{(m)}(\alpha)\neq0.
    \]
\end{enumerate}
\end{methodbox}

\begin{methodbox}[Calculer un PGCD]
Pour calculer \(\pgcd(A,B)\) :
\begin{enumerate}
    \item utiliser l'algorithme d'Euclide ;
    \item rendre le dernier reste non nul unitaire ;
    \item pour trouver les coefficients de Bézout, remonter les divisions.
\end{enumerate}
\end{methodbox}

%=========================================================
\section{Erreurs classiques}

\begin{erreurbox}[Confondre polynôme et fonction]
Un polynôme est un objet formel.  
La fonction associée dépend du corps sur lequel on l'évalue.
\end{erreurbox}

\begin{erreurbox}[Oublier le corps de base]
Le polynôme :
\[
X^2+1
\]
n'a pas de racine dans \(\R\), mais il est scindé dans \(\C\).  
La factorisation dépend du corps.
\end{erreurbox}

\begin{erreurbox}[Croire qu'un degré \(n\) donne toujours \(n\) racines]
Un polynôme de degré \(n\) possède au plus \(n\) racines dans \(\K\).  
Il n'en possède pas forcément \(n\) dans \(\K\).
\end{erreurbox}

\begin{erreurbox}[Mal utiliser le degré du polynôme nul]
On adopte :
\[
\deg(0)=-\infty.
\]
Cette convention évite beaucoup d'exceptions, mais il faut rester prudent dans les raisonnements.
\end{erreurbox}

\begin{erreurbox}[Confondre racine simple et racine multiple]
Une racine \(\alpha\) est multiple si :
\[
P(\alpha)=0
\quad\text{et}\quad
P'(\alpha)=0.
\]
\end{erreurbox}

%=========================================================
\newpage
\section{Exercices progressifs avec corrigés}

\subsection*{Exercice 1 -- Degré et coefficient dominant}

\begin{exercicebox}
Déterminer le degré, le coefficient dominant et dire si le polynôme est unitaire :
\[
P(X)=4X^5-2X^3+X-7.
\]
\end{exercicebox}

\begin{correctionbox}
Le terme de plus haut degré est :
\[
4X^5.
\]
Donc :
\[
\deg(P)=5.
\]
Le coefficient dominant est :
\[
4.
\]
Comme le coefficient dominant n'est pas \(1\), le polynôme n'est pas unitaire.
\end{correctionbox}

\subsection*{Exercice 2 -- Racines évidentes}

\begin{exercicebox}
Montrer que \(1\) est racine de :
\[
P(X)=X^3-2X+1,
\]
puis factoriser \(P\).
\end{exercicebox}

\begin{correctionbox}
On calcule :
\[
P(1)=1-2+1=0.
\]
Donc \(1\) est racine de \(P\), et :
\[
X-1
\]
divise \(P\).

La division donne :
\[
X^3-2X+1=(X-1)(X^2+X-1).
\]
Donc :
\[
P(X)=(X-1)(X^2+X-1).
\]
\end{correctionbox}

\subsection*{Exercice 3 -- Multiplicité}

\begin{exercicebox}
Déterminer la multiplicité de \(1\) comme racine de :
\[
P(X)=X^3-3X^2+3X-1.
\]
\end{exercicebox}

\begin{correctionbox}
On reconnaît :
\[
P(X)=(X-1)^3.
\]
Donc \(1\) est racine de multiplicité \(3\).

On peut aussi vérifier avec les dérivées :
\[
P(1)=0.
\]
\[
P'(X)=3X^2-6X+3,
\qquad
P'(1)=0.
\]
\[
P''(X)=6X-6,
\qquad
P''(1)=0.
\]
\[
P^{(3)}(X)=6,
\qquad
P^{(3)}(1)=6\neq0.
\]
Donc la multiplicité est \(3\).
\end{correctionbox}

\subsection*{Exercice 4 -- PGCD}

\begin{exercicebox}
Calculer le PGCD de :
\[
A=X^3-1,
\qquad
B=X^2-1.
\]
\end{exercicebox}

\begin{correctionbox}
On factorise :
\[
A=X^3-1=(X-1)(X^2+X+1),
\]
\[
B=X^2-1=(X-1)(X+1).
\]
Le facteur commun est :
\[
X-1.
\]
Donc :
\[
\pgcd(A,B)=X-1.
\]
\end{correctionbox}

\subsection*{Exercice 5 -- Bézout}

\begin{exercicebox}
Montrer que :
\[
X
\quad\text{et}\quad
X-1
\]
sont premiers entre eux et trouver une relation de Bézout.
\end{exercicebox}

\begin{correctionbox}
On cherche \(U,V\in\K[X]\) tels que :
\[
UX+V(X-1)=1.
\]
On remarque :
\[
X-(X-1)=1.
\]
Donc :
\[
1\cdot X+(-1)(X-1)=1.
\]
Ainsi :
\[
U=1,
\qquad
V=-1.
\]
Donc \(X\) et \(X-1\) sont premiers entre eux.
\end{correctionbox}

\subsection*{Exercice 6 -- Interpolation}

\begin{exercicebox}
Trouver le polynôme \(P\in\R_2[X]\) tel que :
\[
P(0)=1,\qquad P(1)=2,\qquad P(2)=5.
\]
\end{exercicebox}

\begin{correctionbox}
On cherche :
\[
P(X)=aX^2+bX+c.
\]
La condition \(P(0)=1\) donne :
\[
c=1.
\]
La condition \(P(1)=2\) donne :
\[
a+b+1=2,
\]
donc :
\[
a+b=1.
\]
La condition \(P(2)=5\) donne :
\[
4a+2b+1=5,
\]
donc :
\[
2a+b=2.
\]
En soustrayant :
\[
(2a+b)-(a+b)=2-1,
\]
donc :
\[
a=1.
\]
Puis :
\[
b=0.
\]
Ainsi :
\[
P(X)=X^2+1.
\]
\end{correctionbox}

\subsection*{Exercice 7 -- Polynôme annulateur}

\begin{exercicebox}
Soit \(u\) un endomorphisme tel que :
\[
u^2-3u+2\Id=0.
\]
Montrer que les valeurs propres éventuelles de \(u\) appartiennent à :
\[
\{1,2\}.
\]
\end{exercicebox}

\begin{correctionbox}
Le polynôme :
\[
P(X)=X^2-3X+2
\]
annule \(u\).  
On factorise :
\[
P(X)=(X-1)(X-2).
\]
Si \(\lambda\) est une valeur propre de \(u\), alors :
\[
P(\lambda)=0.
\]
Donc :
\[
(\lambda-1)(\lambda-2)=0.
\]
Ainsi :
\[
\lambda\in\{1,2\}.
\]
\end{correctionbox}

%=========================================================
\newpage
\section{Grand devoir bilan final}

\begin{center}
\Large\textbf{Devoir bilan -- Polynômes}
\end{center}

\subsection*{Exercice 1 -- Questions de cours}

\begin{enumerate}
    \item Définir l'ensemble \(\K[X]\).
    \item Expliquer la différence entre un polynôme et une fonction polynomiale.
    \item Définir le degré d'un polynôme non nul.
    \item Énoncer et démontrer la formule :
    \[
    \deg(PQ)=\deg(P)+\deg(Q).
    \]
    \item Énoncer le théorème de division euclidienne dans \(\K[X]\).
\end{enumerate}

\subsection*{Exercice 2 -- Racines et factorisation}

On considère :
\[
P(X)=X^4-2X^3-X^2+2X.
\]

\begin{enumerate}
    \item Factoriser \(P\) par \(X\).
    \item Montrer que \(1\) et \(-1\) sont racines.
    \item Factoriser complètement \(P\) dans \(\R[X]\).
    \item Donner les multiplicités des racines.
\end{enumerate}

\subsection*{Exercice 3 -- Dérivées et multiplicité}

On considère :
\[
Q(X)=X^4-4X^3+6X^2-4X+1.
\]

\begin{enumerate}
    \item Calculer \(Q(1)\), \(Q'(1)\), \(Q''(1)\), \(Q^{(3)}(1)\), \(Q^{(4)}(1)\).
    \item En déduire la multiplicité de \(1\) comme racine de \(Q\).
    \item Factoriser \(Q\).
\end{enumerate}

\subsection*{Exercice 4 -- PGCD et Bézout}

Soient :
\[
A=X^3-X,
\qquad
B=X^2-1.
\]

\begin{enumerate}
    \item Factoriser \(A\) et \(B\).
    \item Calculer \(\pgcd(A,B)\).
    \item Les polynômes \(A\) et \(B\) sont-ils premiers entre eux ?
    \item Trouver une relation de Bézout entre \(X\) et \(X^2-1\).
\end{enumerate}

\subsection*{Exercice 5 -- Interpolation de Lagrange}

On cherche \(P\in\R_2[X]\) tel que :
\[
P(0)=2,\qquad
P(1)=3,\qquad
P(2)=6.
\]

\begin{enumerate}
    \item Écrire les polynômes de Lagrange associés aux points \(0,1,2\).
    \item Construire \(P\).
    \item Simplifier l'expression de \(P\).
\end{enumerate}

\subsection*{Exercice 6 -- Polynômes annulateurs}

Soit \(u\in\mathcal L(E)\) tel que :
\[
u^3-u=0.
\]

\begin{enumerate}
    \item Donner un polynôme annulateur de \(u\).
    \item Factoriser ce polynôme.
    \item Quelles sont les valeurs propres possibles de \(u\) ?
    \item Que peut-on dire de \(u\) si le corps est \(\R\) ou \(\C\) ?
\end{enumerate}

\newpage

%=========================================================
\section{Correction détaillée du devoir bilan}

\subsection*{Correction de l'exercice 1}

\textbf{1.}  
L'ensemble \(\K[X]\) est l'ensemble des polynômes à coefficients dans \(\K\), c'est-à-dire des expressions formelles :
\[
P(X)=a_0+a_1X+\cdots+a_nX^n,
\]
avec :
\[
a_0,\dots,a_n\in\K.
\]

\textbf{2.}  
Un polynôme est un objet formel défini par une suite finie de coefficients.  
Une fonction polynomiale est l'application obtenue en évaluant le polynôme :
\[
x\mapsto P(x).
\]
Sur un corps infini, deux polynômes qui définissent la même fonction sont égaux.  
Sur un corps fini, ce n'est pas forcément vrai.

\textbf{3.}  
Si :
\[
P(X)=a_0+\cdots+a_nX^n
\]
avec \(a_n\neq0\), alors :
\[
\deg(P)=n.
\]

\textbf{4.}  
Soient \(P,Q\neq0\).  
Écrivons :
\[
P=a_nX^n+\cdots,
\qquad
Q=b_mX^m+\cdots,
\]
avec \(a_n,b_m\neq0\).  
Alors le terme dominant de \(PQ\) est :
\[
a_nb_mX^{n+m}.
\]
Comme \(\K\) est un corps, \(a_nb_m\neq0\).  
Donc :
\[
\deg(PQ)=n+m=\deg(P)+\deg(Q).
\]

\textbf{5.}  
Si \(A,B\in\K[X]\), avec \(B\neq0\), il existe un unique couple \((Q,R)\) tel que :
\[
A=BQ+R
\]
et :
\[
\deg(R)<\deg(B).
\]

\subsection*{Correction de l'exercice 2}

On a :
\[
P(X)=X^4-2X^3-X^2+2X.
\]
On factorise par \(X\) :
\[
P(X)=X(X^3-2X^2-X+2).
\]

On vérifie :
\[
P(1)=1-2-1+2=0,
\]
donc \(1\) est racine.  
Et :
\[
P(-1)=1+2-1-2=0,
\]
donc \(-1\) est racine.

Factorisons le facteur cubique :
\[
X^3-2X^2-X+2.
\]
On regroupe :
\[
X^3-2X^2-X+2
=
X^2(X-2)-1(X-2)
=
(X^2-1)(X-2).
\]
Donc :
\[
X^3-2X^2-X+2
=
(X-1)(X+1)(X-2).
\]
Ainsi :
\[
P(X)=X(X-1)(X+1)(X-2).
\]

Les racines sont :
\[
0,\quad1,\quad-1,\quad2.
\]
Elles sont toutes simples.

\subsection*{Correction de l'exercice 3}

On reconnaît :
\[
Q(X)=X^4-4X^3+6X^2-4X+1=(X-1)^4.
\]
Vérifions par les dérivées.

On a :
\[
Q(1)=1-4+6-4+1=0.
\]

\[
Q'(X)=4X^3-12X^2+12X-4.
\]
Donc :
\[
Q'(1)=4-12+12-4=0.
\]

\[
Q''(X)=12X^2-24X+12.
\]
Donc :
\[
Q''(1)=12-24+12=0.
\]

\[
Q^{(3)}(X)=24X-24.
\]
Donc :
\[
Q^{(3)}(1)=0.
\]

\[
Q^{(4)}(X)=24.
\]
Donc :
\[
Q^{(4)}(1)=24\neq0.
\]

Ainsi \(1\) est racine de multiplicité \(4\).  
Donc :
\[
Q(X)=(X-1)^4.
\]

\subsection*{Correction de l'exercice 4}

On a :
\[
A=X^3-X=X(X^2-1)=X(X-1)(X+1).
\]
Et :
\[
B=X^2-1=(X-1)(X+1).
\]

Donc :
\[
\pgcd(A,B)=X^2-1.
\]
Les polynômes ne sont pas premiers entre eux.

Pour une relation de Bézout entre \(X\) et \(X^2-1\), on remarque :
\[
X^2-(X^2-1)=1.
\]
Donc :
\[
X\cdot X-1\cdot(X^2-1)=1.
\]
Ainsi :
\[
U=X,
\qquad
V=-1
\]
conviennent :
\[
UX+V(X^2-1)=1.
\]

\subsection*{Correction de l'exercice 5}

Les points sont :
\[
a_0=0,\qquad a_1=1,\qquad a_2=2.
\]

Les polynômes de Lagrange sont :

\[
L_0(X)=\frac{X-1}{0-1}\cdot\frac{X-2}{0-2}
=
\frac{(X-1)(X-2)}{2}.
\]

\[
L_1(X)=\frac{X-0}{1-0}\cdot\frac{X-2}{1-2}
=
-X(X-2).
\]

\[
L_2(X)=\frac{X-0}{2-0}\cdot\frac{X-1}{2-1}
=
\frac{X(X-1)}2.
\]

La formule donne :
\[
P(X)=2L_0(X)+3L_1(X)+6L_2(X).
\]
Donc :
\[
P(X)=
2\frac{(X-1)(X-2)}2
+
3[-X(X-2)]
+
6\frac{X(X-1)}2.
\]
Ainsi :
\[
P(X)=(X-1)(X-2)-3X(X-2)+3X(X-1).
\]
Développons :
\[
(X-1)(X-2)=X^2-3X+2,
\]
\[
-3X(X-2)=-3X^2+6X,
\]
\[
3X(X-1)=3X^2-3X.
\]
Donc :
\[
P(X)=X^2-3X+2-3X^2+6X+3X^2-3X.
\]
On obtient :
\[
P(X)=X^2+2.
\]

\subsection*{Correction de l'exercice 6}

On a :
\[
u^3-u=0.
\]
Donc un polynôme annulateur de \(u\) est :
\[
P(X)=X^3-X.
\]
On factorise :
\[
P(X)=X(X^2-1)=X(X-1)(X+1).
\]

Si \(\lambda\) est une valeur propre de \(u\), alors :
\[
P(\lambda)=0.
\]
Donc :
\[
\lambda\in\{-1,0,1\}.
\]

Sur \(\R\) ou \(\C\), le polynôme :
\[
X(X-1)(X+1)
\]
est scindé à racines simples.  
Donc, par le critère usuel de diagonalisation par polynôme annulateur scindé simple, \(u\) est diagonalisable.

%=========================================================
\newpage
\section{Fiche de synthèse finale}

\subsection*{1. Définition}

Un polynôme à coefficients dans \(\K\) est une expression formelle :
\[
P(X)=a_0+a_1X+\cdots+a_nX^n.
\]

\[
\K[X]=\text{ensemble des polynômes à coefficients dans }\K.
\]

Deux polynômes sont égaux si et seulement si leurs coefficients sont égaux.

\subsection*{2. Degré}

Si \(P\neq0\), le degré de \(P\) est le plus grand indice \(k\) tel que le coefficient de \(X^k\) est non nul.

Convention :
\[
\deg(0)=-\infty.
\]

\[
\deg(P+Q)\leq\max(\deg(P),\deg(Q)).
\]

\[
\deg(PQ)=\deg(P)+\deg(Q).
\]

\subsection*{3. Structure}

\[
\K[X]
\]
est un anneau commutatif unitaire et un \(\K\)-espace vectoriel.

Base de \(\K[X]\) :
\[
(1,X,X^2,\dots).
\]

Base de \(\K_n[X]\) :
\[
(1,X,\dots,X^n).
\]

\[
\dim(\K_n[X])=n+1.
\]

\subsection*{4. Division euclidienne}

Si \(A,B\in\K[X]\), avec \(B\neq0\), alors il existe un unique couple \((Q,R)\) tel que :
\[
A=BQ+R,
\qquad
\deg(R)<\deg(B).
\]

\subsection*{5. Racines}

\[
\alpha\text{ racine de }P
\Longleftrightarrow
P(\alpha)=0.
\]

Théorème des racines :
\[
P(\alpha)=0
\Longleftrightarrow
X-\alpha\mid P.
\]

Un polynôme non nul de degré \(n\) possède au plus \(n\) racines distinctes.

\subsection*{6. Multiplicité}

\(\alpha\) est racine de multiplicité \(m\) si :
\[
P(X)=(X-\alpha)^mQ(X),
\qquad
Q(\alpha)\neq0.
\]

Critère par dérivées :
\[
\alpha\text{ racine de multiplicité au moins }m
\]
si et seulement si :
\[
P(\alpha)=P'(\alpha)=\cdots=P^{(m-1)}(\alpha)=0.
\]

Multiplicité exactement \(m\) si, de plus :
\[
P^{(m)}(\alpha)\neq0.
\]

\subsection*{7. Taylor polynomiale}

\[
P(X)=\sum_{k=0}^{\deg(P)}
\frac{P^{(k)}(a)}{k!}(X-a)^k.
\]

\subsection*{8. Factorisation}

Sur \(\C\), tout polynôme non constant est scindé.

Sur \(\R\), tout polynôme non constant se factorise en produit de facteurs irréductibles de degré \(1\) ou \(2\).

\subsection*{9. Irréductibilité}

Dans \(\C[X]\), les irréductibles sont les polynômes de degré \(1\).

Dans \(\R[X]\), les irréductibles sont :
\begin{itemize}
    \item les polynômes de degré \(1\) ;
    \item les polynômes de degré \(2\) sans racine réelle.
\end{itemize}

\subsection*{10. PGCD et Bézout}

Un PGCD \(D\) de \(A\) et \(B\) divise \(A\) et \(B\), et tout diviseur commun de \(A\) et \(B\) divise \(D\).

Bézout :
\[
A\wedge B=1
\Longleftrightarrow
\exists U,V\in\K[X],
\quad
AU+BV=1.
\]

Lemme de Gauss :
\[
A\mid BC,\quad A\wedge B=1
\Longrightarrow
A\mid C.
\]

\subsection*{11. Interpolation de Lagrange}

Pour \(a_0,\dots,a_n\) deux à deux distincts :
\[
L_i(X)=\prod_{j\neq i}\frac{X-a_j}{a_i-a_j}.
\]

\[
L_i(a_j)=\delta_{ij}.
\]

Pour tout \(P\in\K_n[X]\),
\[
P(X)=\sum_{i=0}^n P(a_i)L_i(X).
\]

\subsection*{12. Vandermonde}

\[
\det
\begin{pmatrix}
1&a_0&a_0^2&\cdots&a_0^n\\
1&a_1&a_1^2&\cdots&a_1^n\\
\vdots&\vdots&\vdots&&\vdots\\
1&a_n&a_n^2&\cdots&a_n^n
\end{pmatrix}
=
\prod_{0\leq i<j\leq n}(a_j-a_i).
\]

La matrice est inversible si et seulement si les \(a_i\) sont deux à deux distincts.

\subsection*{13. Polynômes annulateurs}

Si :
\[
P(u)=0,
\]
alors toute valeur propre \(\lambda\) de \(u\) vérifie :
\[
P(\lambda)=0.
\]

Si \(P\) est scindé à racines simples et annule \(u\), alors \(u\) est diagonalisable.

%=========================================================
\section*{Conclusion pédagogique}
\addcontentsline{toc}{section}{Conclusion pédagogique}

Les polynômes forment un chapitre fondamental car ils relient l'algèbre élémentaire, l'arithmétique, l'analyse et l'algèbre linéaire.

Le point central est que \(\K[X]\) se comporte beaucoup comme l'ensemble des entiers :
\begin{itemize}
    \item il existe une division euclidienne ;
    \item on peut définir un PGCD ;
    \item on dispose d'un algorithme d'Euclide ;
    \item on a un théorème de Bézout ;
    \item on peut parler d'irréductibilité.
\end{itemize}

Mais les polynômes possèdent aussi une dimension géométrique et analytique grâce aux racines, aux multiplicités, à la dérivée formelle, à la formule de Taylor et à l'interpolation.

Enfin, ils jouent un rôle essentiel en algèbre linéaire : les polynômes caractéristiques et annulateurs permettent d'étudier les endomorphismes, leurs valeurs propres et leur diagonalisation.

\end{document}